% Copyright (c) 2026 Eugene Joseph M. Ragasa.
% SPDX-License-Identifier: Apache-2.0
% Author: Eugene Joseph M. Ragasa
% Formal proof completed: 2026-10-19
\documentclass[11pt]{article}
\usepackage[T1]{fontenc}
\usepackage{amsmath,amssymb,booktabs}
\usepackage[margin=1in]{geometry}
\usepackage[hidelinks]{hyperref}
\usepackage{tabularx}
\title{Proof Dossier: Ragasa2026A Spectral Field Is Centered}
\author{Eugene Joseph M. Ragasa}
\date{Formal proof completed: 2026-10-19}
\begin{document}
\maketitle
\section{Identity and status}
\begin{tabularx}{\textwidth}{@{}l>{\raggedright\arraybackslash}X@{}}
\toprule
Statement ID & \texttt{BG-RAGASA2026A-THM-005} \\
Status & \texttt{proved} \\
Lean declaration & \path{ProjectKoios.BandGauge.Papers.Ragasa2026A.spectralSplittingField_centered} \\
Module & \path{ProjectKoios/BandGauge/Papers/Ragasa2026A/RankTwoFields.lean} \\
\bottomrule
\end{tabularx}

\section{Exact statement}
Given an exact finite unitary diagonalization
\[
 W_i^\dagger D(q_i)W_i=\operatorname{diag}(d_{i,0},d_{i,1}),
 \qquad d_{i,n}\in\mathbb R,
\]
the spectral deviations satisfy $d_{i,0}+d_{i,1}=0$ for every $i$.

\section{Source and proof}
This theorem \texttt{formalizes} the centered two-state splitting used in manuscript
lines 489--501 and Supplementary Section S1.1 at immutable \path{ksdft2effmass} commit
\begin{center}\small\texttt{99a46cb59d0c6f27ada7d709f194996ef9619d9f}.\end{center}
The centered parent is traceless; unitary conjugation preserves trace; the trace of the
diagonal matrix is the sum of the two supplied real values.

\section{Boundary and trusted base}
The diagonalization is an exact witness, not an eigensolver implementation or floating-
point record. No ordering, continuity, or band-isolation claim is made. Lean and pinned
Mathlib report \texttt{propext}, \texttt{Classical.choice}, and \texttt{Quot.sound}.
No admission or scientific axiom is used. AI-assisted and kernel-checked; no independent
expert review is recorded. Status: \texttt{proved}; working draft.
\end{document}
